\chapter{Integral contractions and the geometric slope inequality}

\begin{lemma}[Tensor coordinates and determinant covariance]
\label{lem:rectangular-tensor-realization}
\uses{lem:rectangular-word-covariance}
\source{bost-specht-manuscript:page-0008,bost-specht-manuscript:page-0009}
\lean{BostConjecture.Geometry.tensorCoordinates,BostConjecture.Geometry.tensorCoordinates_map,BostConjecture.Geometry.map_tensorCoordinates_symm_of_mem_span}
\leanok
Let $R$ be a commutative ring, $M$ a free $R$-module with basis
$(b_i)_{i=1}^a$, and $n\geq0$. The tensor basis identifies $M^{\otimes n}$
with the functions from words $\{1,\ldots,n\}\to\{1,\ldots,a\}$ to $R$.
If $f\in\operatorname{End}_R(M)$ has matrix $U$ in this basis, the
coordinates of $f^{\otimes n}z$ are
\[
 w\longmapsto\sum_v\left(\prod_{q=1}^nU_{w(q),v(q)}\right)z_v.
\]
Consequently, for a tableau $t$ of rectangular shape $(k^a)$ with $k>0$,
the inverse coordinate map sends the $R$-span of the integral orbit
lattice into tensors on which $f^{\otimes n}$ acts by $(\det U)^k$.
\end{lemma}
\begin{proof}
The coordinate of $\bigotimes_q x_q$ at $w$ is
$\prod_q b^*_{w(q)}(x_q)$. Applying $f$ replaces each factor by
$\sum_i U_{w(q),i}b_i^*(x_q)$. Distributing this finite product of
sums gives the stated sum over words. Pure tensors span, so the
formula holds for every tensor. On the scalar span of the rectangular
orbit lattice the word-coordinate action is $(\det U)^k$ by
determinant covariance. Applying the inverse coordinate isomorphism,
which is $R$-linear, proves the asserted tensor identity.
\end{proof}

\begin{lemma}[Change of frame for the rectangular realization]
\label{lem:rectangular-change-frame}
\uses{lem:rectangular-tensor-realization}
\source{bost-specht-manuscript:page-0009}
\lean{BostConjecture.Geometry.tensorCoordinates_symm_change_basis}
\leanok
Let $b,c$ be two bases of the same free rank-$a$ module over a
commutative ring $R$. Write $U$ for the matrix whose columns are the
$c$-coordinates of the vectors of $b$. If $x$ belongs to the scalar
span of a rectangular orbit lattice of shape $(k^a)$ with $k>0$, then
\[
 b^{\otimes n}x=(\det U)^k c^{\otimes n}x.
\]
\end{lemma}
\begin{proof}
The coordinate computation in the preceding lemma applies to a linear
map between modules with different chosen bases: its matrix entries
are the output coordinates of the images of the input basis vectors.
Apply it to the identity map with input basis $b$ and output basis $c$.
The tensor power of the identity is the identity, and the rectangular
word-coordinate action is multiplication by $(\det U)^k$.
\end{proof}

\begin{lemma}[The scalar-extended integral lattice in a free tensor power]
\label{lem:free-integral-tensor-embedding}
\uses{lem:rectangular-tensor-realization,lem:rectangular-change-frame,def:rectangular-specht}
\source{bost-specht-manuscript:page-0009}
\lean{BostConjecture.Geometry.freeLatticeMap,BostConjecture.Geometry.freeLatticeMap_injective,BostConjecture.Geometry.map_freeLatticeMap,BostConjecture.Geometry.freeLatticeMap_change_basis}
\leanok
Let $R$ be a commutative ring, $M$ a free rank-$a$ module with basis $b$,
and $t$ a tableau of rectangular shape $(k^a)$, $k>0$, with integral orbit
lattice $L_t$. There is a linear map
\[
 j_b:R\otimes_{\mathbb Z}L_t\longrightarrow M^{\otimes n},\qquad
 r\otimes x\longmapsto r\,b^{\otimes n}x.
\]
If $R$ is flat over $\mathbb Z$, this map is injective.
For a linear map $f:M\to N$ into another free module with basis $c$,
\[
 f^{\otimes n}j_b(z)=\det([f]_{c,b})^k j_c(z).
\]
In particular $j_b=\det([1]_{c,b})^k j_c$ for two bases of $M$.
\end{lemma}
\begin{proof}
Extend the inclusion $L_t\hookrightarrow\mathbb Z^{\mathrm{words}}$
by scalars to $R$. The scalar extension of the finite coordinate module
is canonically $R^{\mathrm{words}}$; relabel the rectangular alphabet
and apply the inverse tensor coordinate isomorphism. This composite
has the displayed formula. Flatness preserves the injectivity of the
original inclusion, and each remaining map is an isomorphism.
For $z=r\otimes x$, the covariance identity follows from the tensor
coordinate lemma applied to the coefficient extension of $x$; its
coordinates belong to the scalar span by definition. Linearity extends
the identity to every $z$, since elementary tensors span. Taking $f=1$
gives the change-of-frame identity.
\end{proof}

\begin{definition}[Geometric semistability of a Pl\"ucker line]
\label{def:geometric-semistability}
\source{bost-specht-manuscript:page-0009}\notready
For vector spaces $E_1,\ldots,E_N$ over a characteristic-zero field $k$,
put $W=\bigotimes_j E_j$ and $G=\prod_j\SL(E_j)$. An $s$-dimensional
subspace $V\subset W$, $s>0$, has geometrically semistable Pl\"ucker
line if, after extension to an algebraic closure, a homogeneous
$G$-invariant polynomial of positive degree on $\bigwedge^sW$ is
nonzero on a nonzero vector of $\bigwedge^sV$.
\end{definition}

\begin{lemma}[Scalar extension for invariant detection]
\label{lem:invariant-base-change}
\uses{def:geometric-semistability}\notready
The geometric semistability condition is unchanged by extending the
characteristic-zero ground field. In any fixed degree the invariant
polynomials and the kernel of evaluation at a rational vector commute
with field extension.
\end{lemma}
\begin{proof}
For a finite-dimensional representation $M$ of an affine algebraic
group, invariants are the kernel of the linear map
$M\to M\otimes k[G]$, $m\mapsto\rho(m)-m\otimes1$, defined by the
coaction. Tensoring with a field extension is exact, so commutes with
this kernel. Apply this to each homogeneous polynomial space. Evaluation
is a linear map to $k$, and its kernel also commutes with scalar
extension. Its restriction to the invariant space is nonzero after
extension if and only if it was nonzero before extension. Existence of
a positive degree with nonzero restriction is therefore unchanged.
\end{proof}

\begin{lemma}[Determinant-isotypic integral modules]
\label{lem:determinant-isotypic}
\uses{def:rectangular-specht,lem:rectangular-invariants,lem:free-integral-tensor-embedding,lem:degree-calculus}
\source{bost-specht-manuscript:page-0008,bost-specht-manuscript:page-0009}
\notready
For a rank-$a$ projective $\Ocal$-module $A$ and $d=ak>0$, there is
a canonical injection
\[
 \iota_A:\mathcal S_{a,d}(A):=
 (\det A)^{\otimes k}\otimes_\mathbb Z L_{a,d}\longrightarrow A^{\otimes d}.
\]
It is isometric for the determinant and canonical Specht metrics and
satisfies
\[
 \slope(\mathcal S_{a,d}(\overline A))=d\slope(\overline A)-c_a(d).
\]
Its generic fiber is the full $\SL(A_K)$-invariant subspace of
$A_K^{\otimes d}$, with absolutely irreducible isometric $\mathfrak S_d$
action.
\end{lemma}
\begin{proof}
In a free module, a rectangular polytabloid is a permutation of a
tensor product of $k$ unnormalized determinant tensors. Thus every
change-of-frame matrix $g$ acts on it by $(\det g)^k$. For a local
frame $b_1,\ldots,b_a$ the formula
\[
 (b_1\wedge\cdots\wedge b_a)^{\otimes k}\otimes x
 \longmapsto b^{\otimes d}x
\]
is consequently independent of the frame, and glues over $\Ocal$.
The local inclusion of the Specht lattice into the word lattice stays
injective after scalar extension because $\Ocal$ and its localizations
are torsion-free, hence flat, over $\mathbb Z$. This proves injectivity.

At each archimedean place choose an orthonormal frame. Tensor words are
orthonormal, as are the corresponding row tabloids, and the determinant
vector of the frame has norm one. The formula is therefore an isometry.
In particular, no division by a factorial is made: the norm
$\sqrt{a!}$ of an unnormalized determinant tensor already belongs to
the polytabloid's metric.

If $f=\rk_\mathbb Z L_{a,d}$, a basis wedge of
$\Ocal\otimes_\mathbb Z L_{a,d}$ has norm $\covol L_{a,d}$ at every
embedding and generates its determinant module. Its normalized degree
is $-\log\covol L_{a,d}$. The line-twist identity gives the slope
$k\degn(\det A)-f^{-1}\log\covol L_{a,d}$, which is the asserted
formula. The rectangular invariant-space theorem identifies the
generic fiber and its absolute irreducibility. Permutation of tensor
positions preserves the lattice and the ambient Hilbert tensor norm,
and commutes with changes of frame, proving the isometric action.
\end{proof}

\begin{lemma}[Evaluation and permutations in finite projective tensor powers]
\label{lem:projective-tensor-evaluation}
\source{bost-specht-manuscript:page-0010}
\lean{BostConjecture.Tensor.tensorPowerPairing,BostConjecture.Tensor.tensorPowerPairing_tprod_tprod,BostConjecture.Tensor.tensorPowerPairing_reindex_left,BostConjecture.Tensor.tensorPowerPairing_eq_zero_iff}
\leanok
Let $R$ be a commutative ring, $W$ a finite projective $R$-module,
$V$ an $R$-module, $i:V\to W$ linear, and $I$ a finite set.
The canonical finite-projective duality isomorphism defines the linear map
\[
 C_i:V^{\otimes I}\longrightarrow ((W^\vee)^{\otimes I})^\vee,
 \qquad C_i(z)(u)=u(i^{\otimes I}z).
\]
For pure tensors its value is $\prod_{q\in I}f_q(i(v_q))$.
If $\tau$ permutes $I$, then
$C_i(\tau z)(u)=C_i(z)(\tau^{-1}u)$.
Moreover $C_i(z)=0$ if and only if $i^{\otimes I}z=0$.
\end{lemma}
\begin{proof}
The duality isomorphism identifies $\bigotimes_q f_q$ with the functional
$\bigotimes_q w_q\mapsto\prod_q f_q(w_q)$. Its bilinearity and the
linearity of $i^{\otimes I}$ define $C_i$ and give the pure-tensor formula.
Simultaneous permutation of the two pure tensors only permutes the factors
of this product. Bilinearity and spanning by pure tensors extend this
identity to all tensors. The tensor power of $i$ commutes with permutations,
so substitution of $\tau^{-1}u$ gives the asserted dual action identity.
If $i^{\otimes I}z=0$, evaluation vanishes. Conversely the tensor power
of $W$ is projective, hence any nonzero tensor has a linear functional
nonzero on it: embed the projective module as a direct summand of a free
module and choose a nonzero basis coordinate. The finite-projective
duality isomorphism realizes this functional by a tensor of dual vectors,
which contradicts $C_i(z)=0$.
\end{proof}

\begin{lemma}[The canonical contracting pairing]
\label{lem:canonical-contraction}
\uses{lem:determinant-isotypic,lem:projective-tensor-evaluation}
\source{bost-specht-manuscript:page-0010}\notready
Let $i:V\hookrightarrow W=\bigotimes_{j=1}^N E_j$ be a nonzero
submodule with its induced metrics. Set $s=\rk V$, $a_j=\rk E_j$.
For any positive $d$ divisible by $s,a_1,\ldots,a_N$, put
\[
 B_d=\mathcal S_{s,d}(V),\qquad
 D_d=\bigotimes_j\mathcal S_{a_j,d}(E_j^\vee).
\]
Regrouping tensor positions embeds $D_d$ isometrically into
$(W^\vee)^{\otimes d}$. The map
\[
 \Psi_d:B_d\to D_d^\vee,\qquad
 \Psi_d(v)(u)=\langle u,i^{\otimes d}(\iota_V(v))\rangle
\]
is integral, contracting at every archimedean place, and equivariant
for the simultaneous permutation action of $\mathfrak S_d$.
\end{lemma}
\begin{proof}
All maps and pairings in the formula are defined over $\Ocal$, proving
integrality. Each determinant-isotypic embedding is an isometry, and
regrouping Hilbert tensor factors is an isometry. The inclusion $i$ is
isometric by its induced metric, hence so is $i^{\otimes d}$. For
$v\in B_{d,v}$, the functional on the ambient dual tensor space given
by evaluation on $i^{\otimes d}\iota_V(v)$ has norm $\|v\|$.
Its restriction to $D_{d,v}$ has no larger norm. Thus $\|\Psi_d\|\leq1$.
Finally, evaluation satisfies
$\langle u,\tau z\rangle=\langle\tau^{-1}u,z\rangle$, and
$i^{\otimes d}$ commutes with permutations of its $d$ positions.
These two identities give
$\Psi_d(\tau v)(u)=\Psi_d(v)(\tau^{-1}u)$, precisely equivariance
into the dual representation.
\end{proof}

\begin{lemma}[Nonvanishing in unbounded degrees]
\label{lem:contraction-nonzero}
\uses{def:geometric-semistability,lem:invariant-base-change,lem:canonical-contraction,lem:rectangular-invariants}
\source{bost-specht-manuscript:page-0010}\notready
If the Pl\"ucker line of $V_K$ is geometrically semistable, then there
is a positive integer $d_0$, divisible by $s,a_1,\ldots,a_N$, such
that $\Psi_{\ell d_0}\ne0$ for every positive integer $\ell$.
\end{lemma}
\begin{proof}
Extend scalars temporarily to an algebraic closure. Choose
$0\ne\eta\in\det V_K$ and a positive-degree homogeneous invariant
$p$ with $p((\bigwedge^s i)\eta)\ne0$. On $\Hom(V_K,W_K)$ define
$q(A)=p((\bigwedge^s A)\eta)$. It is homogeneous of degree
$d_0=s\deg p>0$, is nonzero at $i$, and is invariant for
$A\mapsto(\bigotimes_jg_j)Ah^{-1}$ with
$(g_1,\ldots,g_N,h)\in\prod_j\SL(E_j)\times\SL(V)$.
A scalar $a_j$-th root of unity in the $j$th factor acts by scalar
multiplication on $A$. Since $q$ is nonzero and homogeneous, invariance
forces its $d_0$th power to be one for every such root, so $a_j\mid d_0$.
Also $s\mid d_0$ by its definition.

For $d=\ell d_0$, polarization identifies $q^\ell$ with a symmetric
tensor in $(\Hom(V_K,W_K)^\vee)^{\otimes d}$: the inverse to the
symmetric quotient is averaging over permutations with factor $1/d!$.
This identification commutes with the group action and evaluation on
$A^{\otimes d}$. After regrouping, its invariant tensor belongs to
\[
 \left(\bigotimes_j(E_j^\vee)^{\otimes d}\otimes V^{\otimes d}\right)
 ^{\prod_j\SL(E_j)\times\SL(V)}=D_{d,\overline K}\otimes B_{d,\overline K}.
\]
For independent group actions, the equality of invariants follows by
expanding in a basis of every factor except one and imposing invariance
in that factor successively; the rectangular invariant-space theorem
identifies each resulting factor. Write the tensor as
$\sum_\nu u_\nu\otimes v_\nu$. Evaluation gives
$q(i)^\ell=\sum_\nu\Psi_d(v_\nu)(u_\nu)\ne0$.
Therefore $\Psi_d$ is nonzero after extension, and hence before it.
The polarized tensor is used only for nonvanishing; no integral or
metric bound for its coefficients enters the argument.
\end{proof}

\begin{theorem}[Geometric slope inequality]
\label{thm:geometric-slope}
\uses{lem:contraction-nonzero,lem:canonical-contraction,lem:determinant-isotypic,lem:irreducible-semistable,lem:degree-calculus,thm:specht-sublinear-covolume}
\source{bost-specht-manuscript:page-0009,bost-specht-manuscript:page-0010}
\notready
For arbitrary positive-rank Hermitian bundles $E_1,\ldots,E_N$ and
a nonzero submodule $V\subseteq\bigotimes_j E_j$, geometric
semistability of the Pl\"ucker line of $V_K$ implies
\[
 \slope(\overline V)\leq\sum_j\slope(\overline E_j).
\]
The factors need not be arithmetically semistable.
\end{theorem}
\begin{proof}
Take the unbounded sequence $d=\ell d_0$ from the nonvanishing lemma.
The kernel of $\Psi_{d,K}$ is invariant for the simultaneous
$\mathfrak S_d$ action. Absolute irreducibility of $B_{d,K}$ and
nonvanishing force that kernel to be zero. Torsion-freeness gives
injectivity over $\Ocal$. Let $J_d$ be the actual image, with metric
induced from $D_d^\vee$. The isomorphism of modules $B_d\to J_d$
is contracting, so $\slope(B_d)\leq\slope(J_d)$.

Each factor of $D_d$ has absolutely irreducible isometric permutation
action. Their tensor product is absolutely irreducible for the
independent product of permutation groups. Its dual is semistable by
the irreducible-action lemma. Thus, including when $J_d$ is not
saturated, $\slope(J_d)\leq\slope(D_d^\vee)$. The determinant and
average tensor-slope formulas give
\[
 d\slope(V)-c_s(d)\leq\slope(J_d)
 \leq d\sum_j\slope(E_j)+\sum_jc_{a_j}(d).
\]
Divide by $d$ and let $\ell\to\infty$. The sublinear Specht covolume
theorem makes every error term tend to zero, giving the claim.
\end{proof}
